\documentclass[12pt,a4paper]{article}

% Υποστήριξη XeLaTeX & Γραμματοσειρές
\usepackage{fontspec}
\usepackage{xunicode}
\usepackage{xltxtra}

% Ορισμός Liberation Serif για Ελληνικά και Αγγλικά
\setmainfont{Liberation Serif}

% Μαθηματικά πακέτα
\usepackage{amsmath,amssymb,amsfonts,mathtools}

% Διαμόρφωση σελίδας
\usepackage{geometry}
\geometry{top=2.5cm, bottom=2.5cm, left=2.5cm, right=2.5cm}

% Χρώματα και διακόσμηση
\usepackage{xcolor}
\usepackage{enumitem}

% Προσαρμογή επικεφαλίδων
\usepackage{titlesec}
\usepackage{hyperref}

\titleformat{\section}{\Large\bfseries\color{black}}{\thesection}{1em}{}
\titleformat{\subsection}{\normalfont\bfseries\color{black}}{\thesubsection}{1em}{}

\usepackage{tikz}
\usepackage{fancyhdr}
\usepackage{caption}
\pagestyle{fancy}
\renewcommand{\headrulewidth}{0pt}
\fancyhf{} 
\fancyfoot[L]{\footnotesize Theolaos}
\fancyfoot[R]{\small \thepage}

\usetikzlibrary{decorations.pathreplacing}
\usetikzlibrary{arrows.meta}

\setlist[itemize]{left=1.5em, label=\textbullet}
\setlength{\parskip}{0pt}
\setlength{\parindent}{0pt}

\begin{document}

\begin{center}
    \LARGE \textbf{Μεθοδολογία Διαφορικών Εξισώσεων 1$^{\text{ου}}$ Βαθμού}\\[0.5ex]
    \large Σημειώσεις Μαθηματικά ΙΙ \\
    \normalsize Βρες και άλλες σημειώσεις: \url{https://github.com/theolaos/simeiwseis}
\end{center}

\section*{Μεθοδολογία Χωριζομένων Μεταβλητών (ΚΑΤ. 1)}
\vspace{-1em}
\hrule
\vspace{1em}


\[
\frac{dy}{dx} = f(x,y) \iff \frac{dy}{dx} = f(x)f(y)
\]

Σκοπός μας είναι να χωρίσουμε το 2ο μέλος σε γινόμενο δύο συναρτήσεων: μία που περιέχει μόνο το $x$ και μία που περιέχει μόνο το $y$.

Έπειτα διαχωρίζουμε τις μεταβλητές στα δύο μέλη:
\[
\frac{dy}{f(y)} = f(x)\,dx
\]

Ή ισοδύναμα, προσπαθούμε εξ αρχής να φέρουμε την εξίσωση στη μορφή:
\[
f(y)\,dy = f(x)\,dx
\]

Στη συνέχεια ολοκληρώνουμε και τα δύο μέλη και λύνουμε ως προς $y$.

\textbf{Προσοχή:} Μην ξεχνάς τη σταθερά $+C$ στο τέλος!


\section*{Μεθοδολογία Ομογενών Συναρτήσεων (ΚΑΤ. 2)}
\vspace{-1em}
\hrule
\vspace{1em}

Φέρνουμε τη ΔΕ στη μορφή $\frac{dy}{dx} = f(x,y)$ προκειμένου να ελέγξουμε αν η $f(x,y)$ είναι ομογενής.

Μια συνάρτηση είναι ομογενής $n$-βαθμού όταν ισχύει:
\[
f(\lambda x, \lambda y) = \lambda^n f(x,y)
\]

Εφόσον ισχύει η παραπάνω ιδιότητα, ακολουθούμε τα εξής βήματα:
\begin{enumerate}[leftmargin=*]
    \item Θέτουμε $u = \frac{y}{x} \implies y = u \cdot x$.
    \item Παραγωγίζουμε ως προς $x$:
    \[
    \frac{dy}{dx} = \frac{du}{dx} x + u
    \]
    \item Αντικαθιστούμε στη αρχική ΔΕ:
    \[
    \frac{du}{dx} x + u = f(1,u)
    \]
    \item Επιλύουμε τη νέα ΔΕ ως προς $u$ με τη \textbf{μέθοδο χωριζομένων μεταβλητών (ΚΑΤ. 1)}.
\end{enumerate}

\section*{Γραμμική Μεθοδολογία (ΚΑΤ. 3)}
\vspace{-1em}
\hrule
\vspace{1em}

Εφαρμόζεται όταν η ΔΕ έχει τη μορφή:
\[
\frac{dy}{dx} + P(x)y = Q(x)
\]

Η γενική λύση δίνεται απευθείας από τον τύπο:
\[
y = e^{-\int P(x)dx} \left( C + \int Q(x) e^{\int P(x)dx} \, dx \right)
\]

\textit{Σημείωση:} Απαιτείται μόνο ο υπολογισμός δύο ολοκληρωμάτων.

\section*{Μεθοδολογία Τύπου Bernoulli (ΚΑΤ. 4)}
\vspace{-1em}
\hrule
\vspace{1em}

Η διαφορική εξίσωση Bernoulli έχει τη μορφή:
\[
\frac{dy}{dx} + P(x)y = Q(x)y^n
\]

\textbf{Ειδικές περιπτώσεις:}
\begin{itemize}
    \item Αν $n = 0 \implies$ Γραμμική ΔΕ (ΚΑΤ. 3).
    \item Αν $n = 1 \implies$ Χωριζομένων Μεταβλητών (ΚΑΤ. 1).
\end{itemize}

\textbf{Γενική περίπτωση ($n \neq 0, 1$):}
\begin{enumerate}[leftmargin=*]
    \item Πολλαπλασιάζουμε τα δύο μέλη με $y^{-n}$:
    \begin{equation}
    y^{-n} \frac{dy}{dx} + P(x) y^{1-n} = Q(x)
    \end{equation}
    
    \item Θέτουμε $z = y^{1-n}$ και παραγωγίζουμε ως προς $x$:
    \[
    \frac{dz}{dx} = (1-n) y^{-n} \frac{dy}{dx} \iff \frac{1}{1-n} \frac{dz}{dx} = y^{-n} \frac{dy}{dx}
    \]
    
    \item Αντικαθιστούμε στην εξίσωση (1):
    \[
    \frac{dz}{dx} + (1-n)P(x)z = (1-n)Q(x)
    \]
    \item Η εξίσωση είναι πλέον \textbf{γραμμική} ως προς $z$ και λύνεται με τη μεθοδολογία ΚΑΤ. 3.
\end{enumerate}

\section*{Μεθοδολογία Τύπου Clairaut (ΚΑΤ. 5)}
\vspace{-1em}
\hrule
\vspace{1em}

Η ΔΕ Clairaut έχει τη μορφή:
\[
y = x y' + \phi(y')
\]

Θέτοντας $p = y'$, η εξίσωση παίρνει τη μορφή $\left( x + \phi'(p) \right) p' = 0$.

\subsection*{Λύσεις της Clairaut}
\begin{itemize}
    \item \textbf{Γενική Λύση:} Προκύπτει για $p' = 0 \implies p = C$, οπότε:
    \[
    y = C x + \phi(C)
    \]
    \item \textbf{Ιδιάζουσες Λύσεις (Περιβάλλουσα):} Προκύπτουν από το σύστημα:
    \[
    x + \phi'(p) = 0 \quad \text{και} \quad y = x p + \phi(p)
    \]
\end{itemize}

\subsection*{Βήματα Επίλυσης}
\begin{enumerate}[leftmargin=*]
    \item[0.] Φέρνουμε τη ΔΕ στη μορφή Clairaut.
    \item Αναγνωρίζουμε τη συνάρτηση $\phi(x)$.
    \item Γράφουμε τη Γενική Λύση: $y = C x + \phi(C)$.
    \item Βρίσκουμε τις Ιδιάζουσες Λύσεις απαλείφοντας την παράμετρο $p$.
\end{enumerate}

\subsection*{Παράδειγμα (Διαφάνεια 29, Άσκηση 3)}

Δίνεται η εξίσωση:
\[
y(y')^2 + 2xy' = y \iff (y y')^2 + 2x(y y') = y^2
\]

Θέτουμε $u = y^2 \implies u' = 2y y' \implies y y' = \frac{1}{2} u'$.

Αντικαθιστώντας, παίρνουμε:
\[
\left( \frac{1}{2} u' \right)^2 + x u' = u \iff u = x u' + \left( \frac{u'}{2} \right)^2
\]
Η οποία είναι μορφής Clairaut ως προς $u$, με $\phi(t) = \left(\frac{t}{2}\right)^2$.

\begin{itemize}
    \item \textbf{Γενική Λύση ως προς $u$:}
    \[
    u = C x + \left(\frac{C}{2}\right)^2 \implies y^2 = C x + \left(\frac{C}{2}\right)^2
    \]
    
    \item \textbf{Ιδιάζουσα Λύση:}
    \[
    x + \phi'(u') = 0 \iff x + \frac{u'}{2} = 0 \implies u' = -2x
    \]
    Αντικαθιστώντας το $u' = -2x$ στη ΔΕ:
    \[
    u = x(-2x) + \left(\frac{-2x}{2}\right)^2 = -2x^2 + x^2 = -x^2 \implies y^2 = -x^2
    \]
    (\textit{Αδύνατη στο πεδίο των πραγματικών αριθμών}).
\end{itemize}

\section*{Μεθοδολογία Τύπου Lagrange (ΚΑΤ. 6)}
\vspace{-1em}
\hrule
\vspace{1em}


Η ΔΕ του Lagrange έχει τη γενική μορφή:
\[
y = x \phi(y') + f(y')
\]

\end{document}